% Build the layer, Chapter 2, fifth build: who may revoke a card.
% Inserted after the fourth build (\S\ref{sec:anchor-build-told}), whose
% gate it fixes. It closes the limitation that the third and fourth builds
% each recorded, and works through the two exercises that ask for it.
%
% No code or output in this file is typed by hand. \pdexcerpt, \pdrun and
% \pdrunlines print files under excerpts/, which the kernel's
% tools/excerpts.mjs cuts from its source and its recorded runs (see
% figures/pd-listing.tex). To change a listing, change the kernel and cut
% again. The prose names functions, never line numbers, so that a listing
% which moves in its file does not make a sentence wrong.
%
% \pdlistingroom{n} before a listing of n lines keeps it on one page. When a
% listing changes length, change the number with it.
\subsection{Build the layer: who may revoke a card}
\label{sec:anchor-build-warranted}

\noindent\emph{Express lane: the last two gates revoke a card for
anyone who presents its chain, and a chain is not a secret. Requiring an
answer to a challenge keeps an outsider out at one gate, and leaves that
gate with nothing the others can verify. The gate that works revokes a card
only on a signature over the harbor's key and the card's id, made by the
harbor or by a holder the chain names. That signature answers no challenge,
so every gate can verify the same one.}

Two limitations have been waiting. \S\ref{sec:anchor-build-taken-back} said
that its gate does not check who is revoking a card.
\S\ref{sec:anchor-build-told} made that matter more. The notice of a
revocation is the card's chain, so every gate that receives the notice
holds the ancestors of that card, and so does anyone reading the network
traffic. This section builds
the rule both sections asked for, in \pdkernelname{}.
Exercises~\ref{ex:anchor-taken-who} and~\ref{ex:anchor-told-above} ask the
reader to design it first, and they are worth trying before reading on.

\subsubsection{Break: an outsider revokes a card}

Alice holds a card from the harbor. She has delegated to Heron, and he
has delegated to Finch. Wren holds none of the three cards. She has a copy
of Finch's chain, copied as he sent it to the gate, and since
\S\ref{sec:anchor-build-holder} a copy grants her nothing. The gate is
the fourth build's, of Listing~\ref{alg:anchor-told-tell}.

\pdlistingroom{16}
\pdrunlines[caption={An outsider with a copy of a chain has its last card
revoked.},label={alg:anchor-warranted-break}]{outsider-verified}{1,11-25}{1 and 11\textendash 25}

Wren gains nothing she can use, and Finch loses his card at the request of
someone who holds no card at all. When Heron comes a minute later to revoke
the card himself, the gate answers that it is revoked, as it already was.
This gate also gossips. The chain is its notice, and each gate that
receives it revokes Finch's card in turn.

\subsubsection{A first repair: require a proof, as a presentation does}

\S\ref{sec:anchor-build-holder} met this problem with presentations, where
a copied card was admitted, and its fix is available. Let the gate issue a
challenge, and revoke a card only for an answer signed with a key
that has authority over the card: the harbor's, or one that a card of the
chain names. This is the design that the solution to
Exercise~\ref{ex:anchor-taken-who} arrives at.

\pdlistingroom{29}
\pdexcerpt[caption={The first attempt. A card is revoked for whoever
answers this gate's challenge.},label={alg:anchor-warranted-challenged}]{ch02e-challenged}

Run the break against this gate and it is fixed. Wren is refused when she
sends no answer. She is refused again with each answer she can make: one
signed with her own key, which no card of the chain names, and one that
names Heron's key over a signature of her own. Then Heron answers
a challenge, and Finch's card is revoked for him.

\subsubsection{Break again: the answer cannot be verified at another gate}

The fourth build asked what one gate can send another. Ask it of this
gate. Two gates trust one harbor, as in \S\ref{sec:anchor-build-told}, and
Alice revokes Heron's card at the near one.

\pdlistingroom{20}
\pdrunlines[caption={A card revoked on an answer at one gate. The gates
exchange notices, and it is still accepted at the
other.},label={alg:anchor-warranted-unheard}]{two-gates-challenged}{1,11-29}{1 and 11\textendash 29}

The line at $t=240000$ is where it goes wrong: the gates exchange notices,
and neither adds anything. The near gate sent what it had, which is
the chain and Alice's answer. The far gate checked them as it checks
anything sent to it, and refused them as \texttt{unknown-challenge},
because the answer is to a challenge the far gate never issued. It could
skip that check for a notice from another gate. Then it would be trusting
what a gate says, which is the unverified notice of
Listing~\ref{alg:anchor-told-rumour} again.
So the first repair keeps the outsider out and breaks the last build. Wren
goes to the other gate, and 60 seconds after the card was revoked she
is admitted there.

\subsubsection{Fix: a signature with no challenge in it}

Why does a presentation need a challenge at all? The same presentation
accepted twice admits someone twice, and the challenge is what makes the
second one invalid. Revoking a card is different. Revoked a second time,
the card is in the state the first revocation left it: the list has one
entry for it, with the same expiry. A second copy of the request changes
nothing, so nothing has to make it invalid.

So remove the challenge. What is left is a signature that says which card,
at which harbor. This book calls it a \pdterm{signed revocation}.

\pdlistingroom{21}
\pdexcerpt[caption={Signing a revocation. The text contains no
challenge.},label={alg:anchor-warranted-sign}]{ch02e-sign}

Listing~\ref{alg:anchor-warranted-text} is the same text in the core.

\pdlistingroom{8}
\pdexcerpt[caption={What a revocation is a signature
over.},label={alg:anchor-warranted-text}]{ch02e-text}

A revocation carries no card id. The core computes the text from the gate's
own key and the last card of the chain it was sent, as it computes the text
of a proof (Listing~\ref{alg:anchor-holder-text}). So a revocation contains
nothing that a gate has to trust, and sent with another chain
it is a signature over the wrong text.

The core checks a revocation after it has checked the chain.

\pdlistingroom{29}
\pdexcerpt[caption={The core's check of a signed
revocation.},label={alg:anchor-warranted-core}]{ch02e-authorized}

It checks two things, in this order. First, whose key signed? It must be
the harbor's, or a key that some card of the chain names. Those are the keys
of the card's own holder and of the holder of every ancestor. The holder of
a descendant has no authority over the card, and neither has the holder of
a sibling, whose card is in another chain. Second, is it a signature over
the right text?

The gate keeps the revocation with the chain, and the two together are its
notice.

\pdlistingroom{36}
\pdexcerpt[caption={The gate that revokes a card on a signed revocation,
and forwards the revocation.},label={alg:anchor-warranted-gate}]{ch02e-gate}

\texttt{receive} passes what it is sent to \texttt{revoke}, as the fourth
build's did. There is still one way onto the list, and now it checks who is
asking.

% A page may break here. See the note in anchor-build-the-verifier.tex.
\pagebreak[0]
\keyidea[What may be repeated needs no challenge]{A challenge makes a
message valid once. Before adding one, ask what a second copy of the
message could do. If it could do nothing, the message can stay valid, and
then any gate can verify it and anyone can forward it.}

Two published designs each have half of this rule. OAuth lets a token be
revoked by the client it was issued to, which is the holder revoking its own
card~\pdcite{rfc7009revocation}. UCAN lets a delegation be revoked by
any issuer in its chain, which is the holders of its ancestors, and its
record leaves the nonce empty because revoking twice changes
nothing~\pdcite{ucan2026revocation}. \pdkernelname{} takes both halves.

\subsubsection{Verify: the same runs, and every signer}

The outsider first (Listing~\ref{alg:anchor-warranted-refused}).

\pdlistingroom{26}
\pdrunlines[caption={The outsider against the fix. She is refused each way
she tries, and the holder of the parent card is
not.},label={alg:anchor-warranted-refused}]{outsider-core}{1,13-37}{1 and 13\textendash 37}

Each refusal has its own reason, and each reason is a line of
Listing~\ref{alg:anchor-warranted-core}. With no revocation there is nothing
to check. Of the two she signs, one is signed with a key the chain does not
name, and the other names Heron's key over a signature that is not his.

Now a valid revocation, sent with a card it was not signed for. Heron has
delegated to Wren as well, a sibling of Finch's card, and revokes it. Wren
keeps a copy of his revocation.

\pdlistingroom{22}
\pdrunlines[caption={A revocation of one card, sent again with two
others.},label={alg:anchor-warranted-misdirected}]{wrong-card-core}{1,13-33}{1 and 13\textendash 33}

Heron's key has authority over all three cards, so the first check does not
stop her. The second does: the text the core computes for Finch's card is
not the text Heron signed.

The two gates of \S\ref{sec:anchor-build-told} need no new listing. Under
this gate, \texttt{harbor scenario two-gates --design core} prints
Listing~\ref{alg:anchor-told-heard} again, line for line, and the
false notice's run prints Listing~\ref{alg:anchor-told-refused}. In that run
Wren holds Heron's key, so she can sign a revocation of his card, and she
sends one with her changed chain. The core refuses the chain before it
reads the revocation.

A scenario tries a few signers. The rule is small enough to try them all.
Take the chain of the first run, and let five signers each sign a
revocation of each of its cards. Four are the harbor and the chain's
holders, and the fifth is Wren.

\pdlistingroom{9}
\pdrunlines[caption={Every signer, and every card of a chain of
three.},label={alg:anchor-warranted-every}]{revokers}{1,5-12}{1 and 5\textendash 12}

The table is a staircase. Each holder can have its own card revoked and
every card delegated from it, and none of its ancestors. The harbor can have
any card revoked, and Wren none. Figure~\ref{fig:anchor-warrant-who} draws
it.

% Brief (grammar: every case in a small matrix; atlas
%   V/fig:anchor-warrant-who).
% Reader question: whose signed revocation revokes which card of a chain?
% Claim: a card is revoked for the harbor, and for whoever holds it or one
% of its ancestors. The revocations the rule refuses are the ones a core
% built without the check of the signer accepts.
% Evidence: 15 revocations, five signers for each of three cards, sent to
% the core's gate and to a core built without the check of the signer
% [internal, harbor revokers; sessions revokers.txt and
% revokers-any-signer.txt, and
% tests/harbor-research/test_anchor_warranted_numbers.py checks every mark
% against them].
% Must distinguish: a revocation the rule accepts from one accepted only
% with the guard removed; a holder's own card from its ancestors and its
% descendants.
% Grammar chosen: signers down and cards across, both in the order of the
% chain, one mark for each revocation, so that the rule shows as a
% staircase.
% Grammar rejected: the count alone, which does not say which nine; the
% chain drawn as a tree with an arrow for each revocation, where fifteen
% arrows cross.
% Five-second test: what shape do the filled marks make?
\begin{figure}[H]
\centering
\pdfigurehue{pdcobalt}%
\begin{tikzpicture}[pd figure]
  \pgfplotsset{
    revoked/.style={only marks,mark=*,mark size=3.2pt,draw=pdfocus,
      mark options={fill=pdfocus,draw=pdpage,line width=.3pt}},
    unchecked/.style={only marks,mark=o,mark size=3.2pt,draw=pdinkmuted,
      mark options={draw=pdinkmuted,line width=.6pt}},
  }
  \begin{axis}[pd axis,width=5.4cm,height=4.4cm,
    xmin=0.5,xmax=3.5,ymin=0.4,ymax=5.6,
    xtick={1,2,3},xticklabels={Alice's,Heron's,Finch's},
    ytick={1,2,3,4,5},yticklabels={Wren,Finch,Heron,Alice,the harbor},
    xlabel={the card the revocation names},
    ylabel={who signs it}]
    % The staircase: below it and left of it, the rule refuses.
    \addplot[pd guide,mark=none] coordinates {(0.5,3.5) (1.5,3.5) (1.5,2.5) (2.5,2.5) (2.5,1.5) (3.5,1.5)};
    % series: revoked
    \addplot[revoked] coordinates {(1,5) (2,5) (3,5) (1,4) (2,4) (3,4) (2,3) (3,3) (3,2)};
    % series: revoked only without the check of the signer
    \addplot[unchecked] coordinates {(1,3) (1,2) (2,2) (1,1) (2,1) (3,1)};
  \end{axis}
\end{tikzpicture}
\caption{\textbf{A card is revoked for the harbor, and for whoever holds
it or one of its ancestors.} Filled marks are the nine revocations the gate
accepts. Open marks are the six it refuses, which a core built without the
check of the signer accepts too. [internal, \texttt{harbor revokers}]}
\label{fig:anchor-warrant-who}
\end{figure}


Each check this build added is a guard that can be removed. Build the
core without the check of whose key signed
(Listing~\ref{alg:anchor-warranted-any-signer}).

\pdlistingroom{9}
\pdrunlines[caption={The fix with the check of the signer removed. Every
revocation is accepted.},label={alg:anchor-warranted-any-signer}]{revokers-any-signer}{1,5-12}{1 and 5\textendash 12}

Six cells change, and they are the six the rule refused. Then build it
without the check of the signature
(Listing~\ref{alg:anchor-warranted-unwarranted}).

\pdlistingroom{14}
\pdrunlines[caption={The fix with the signature check removed. One
revocation revokes three
cards.},label={alg:anchor-warranted-unwarranted}]{wrong-card-any-revocation}{1,18-30}{1 and 18\textendash 30}

The third guard is the refusal of a request that carries no revocation.
Without it the outsider's first attempt in
Listing~\ref{alg:anchor-warranted-refused} is accepted, and the run
ends as Listing~\ref{alg:anchor-warranted-break} did.

The check is also three ProVerif models
(Listing~\ref{alg:anchor-warranted-models}).

\pdlistingroom{16}
\pdrun[caption={ProVerif~2.05 on models of a signed
revocation.},label={alg:anchor-warranted-models}]{models-revokers}

\texttt{Signed} marks a key signing a revocation of a card at a harbor, and
\texttt{Revoked} marks a gate revoking that card on it.
\texttt{HeronsRevokedBy} marks the gate revoking Heron's card on a
revocation signed with the key named. In every model
the harbors and holders sign whatever the attacker asks them to. Under
\texttt{authorized.pv}, which is the fix, a gate still revokes a card
only on a revocation that the key it accepted signed, for that card at that
harbor. It never revokes Heron's card on the key of Finch, who holds
a card delegated from Heron's, or on Wren's. In \texttt{any-signer.pv} every
line is false. \texttt{no-harbor-in-the-text.pv} leaves the harbor's key out
of the text a revocation signs, and a revocation signed for one harbor is
accepted at another (Exercise~\ref{ex:anchor-warranted-harbor}).

Two lines of the fix say \texttt{false}, and both are meant to. The last
shows that a card can be revoked at all. The \texttt{inj-event} line asks
whether a revocation is accepted no more often than it was signed, which is
what \S\ref{sec:anchor-build-holder} required of a proof. A revocation is
accepted as often as it is sent. That is the property this section set out
to get.
\pdscope{The runs cover one chain of three cards and five signers, which
is 15 revocations, at gates in one process. The models cover the same chain
with any number of revocations and an attacker who chooses what is signed.
Their cryptography cannot be broken. Longer chains depend on the comparison
of keys being the same for every card.}

\begin{pdboundary}
\setlength{\parskip}{4pt plus 1pt}
A revocation is verified with a key, and a key belongs to whoever has it.
Wren with Heron's key can revoke his card and every card delegated from it.
She can revoke none of its ancestors, and she could already do whatever
Heron could.

A signed revocation cannot itself be revoked. Once it has been seen, it
works at every gate of its harbor until the card expires, for whoever sends
it. A holder who revokes a card by mistake signs a new card, and the old one
stays revoked (Design invariant~\ref{prop:no-reanim}).

A revocation shows whose key signed, and nothing about who sent it or when
it was signed. A gate that needs those facts has to record them itself.

The notice is still a chain, and now a revocation with it, so a gate that
receives one learns who revoked the card as well as who held which rights,
delegated from whom.
\end{pdboundary}
